\section{微积分在几何与物理中的应用}\label{sec:8-2}

\subsection{直角坐标下平面图形的面积}\label{subsec:8-2-1}

求由 $x = a, x = b, y = f(x), y = g(x)$（其中 $g(x) \leqslant f(x)$）围成的面积（如图~\ref{fig:8-2}），显然面积
\[
A = \int_a^b [f(x) - g(x)]\,\mathrm{d}x.
\]
公式并不要求 $g(x) \geqslant 0$，面积微元
\[
\mathrm{d}A = [f(x) - g(x)]\,\mathrm{d}x.
\]
如果 $g(x) = 0$，$f(x) \geqslant 0$，则图形为曲边梯形，面积微元为
\[
\mathrm{d}A = f(x)\,\mathrm{d}x.
\]
如果图形如图~\ref{fig:8-3} 所示，则由 $y = \alpha, y = \beta, x = \varphi(y), x = \psi(y)$（其中 $\psi(y) \leqslant \varphi(y)$）所围成的面积为
\[
A = \int_\alpha^\beta [\varphi(y) - \psi(y)]\,\mathrm{d}y,
\]
面积微元
\[
\mathrm{d}A = [\varphi(y) - \psi(y)]\,\mathrm{d}y.
\]

\begin{figure}[htbp]
\centering
\begin{minipage}[t]{0.45\textwidth}
\centering
\includegraphics[width=0.9\textwidth]{fig8-2.png}
\figcap{图 8-2}\label{fig:8-2}
\end{minipage}
\hfill
\begin{minipage}[t]{0.45\textwidth}
\centering
\includegraphics[width=0.9\textwidth]{fig8-3.png}
\figcap{图 8-3}\label{fig:8-3}
\end{minipage}
\end{figure}

\begin{example}
求 $y = 2, y = x, xy = 1$ 所围图形的面积（见图~\ref{fig:8-4}）．
\end{example}

\begin{figure}[htbp]
\centering
\includegraphics[width=0.45\textwidth]{fig8-4.png}
\figcap{图 8-4}\label{fig:8-4}
\end{figure}

\begin{proof}[解]
先求出三条曲线的交点为 $\left(\frac{1}{2}, 2\right)$，$(1, 1)$，$(2, 2)$，则面积
\begin{align*}
A &= \int_{1/2}^1 \left(2 - \frac{1}{x}\right)\mathrm{d}x + \int_1^2 (2 - x)\,\mathrm{d}x \\
&= (2x - \ln x)\Big|_{1/2}^1 + \left(2x - \frac{x^2}{2}\right)\Bigg|_1^2 \\
&= \frac{3}{2} - \ln 2.
\end{align*}
也可以用另一个公式计算
\[
A = \int_1^2 \left(y - \frac{1}{y}\right)\mathrm{d}y = \left(\frac{y^2}{2} - \ln y\right)\Bigg|_1^2 = \frac{3}{2} - \ln 2.
\]
\end{proof}

\begin{example}
求由参数方程
\[
\begin{cases}
x = a + a\cos t, \\
y = a + a\sin t,
\end{cases}
\quad 0 \leqslant t \leqslant 2\pi
\]
所围图形的面积（见图~\ref{fig:8-5}）．
\end{example}

\begin{figure}[htbp]
\centering
\includegraphics[width=0.4\textwidth]{fig8-5.png}
\figcap{图 8-5}\label{fig:8-5}
\end{figure}

\begin{proof}[解]
记上半圆周曲线为 $y_2(x)$，下半圆周曲线为 $y_1(x)$，则
\[
A = \int_0^{2a} [y_2(x) - y_1(x)]\,\mathrm{d}x.
\]
作变量代换，则
\begin{align*}
A &= \int_\pi^0 a(1 + \sin t)\,\mathrm{d}[a(1 + \cos t)] - \int_\pi^{2\pi} a(1 + \sin t)\,\mathrm{d}[a(1 + \cos t)] \\
&= -\int_0^{2\pi} a(1 + \sin t)\,\mathrm{d}[a(1 + \cos t)] \\
&= \int_0^{2\pi} a^2(1 + \sin t)\sin t\,\mathrm{d}t \\
&= a^2 \int_0^{2\pi} \sin^2 t\,\mathrm{d}t = \frac{a^2}{2}\int_0^{2\pi} (1 - \cos 2t)\,\mathrm{d}t \\
&= \frac{a^2}{2}\cdot 2\pi - \frac{a^2}{4}\sin 2t\Big|_0^{2\pi} = \pi a^2.
\end{align*}
\end{proof}

\begin{example}
求椭圆 $\frac{x^2}{a^2} + \frac{y^2}{b^2} = 1$ 所围图形的面积．由对称性，只需求第一象限内的面积（图~\ref{fig:8-6}）．
\end{example}

\begin{proof}[解]
\[
A = 4\int_0^a \frac{b}{a}\sqrt{a^2 - x^2}\,\mathrm{d}x = \frac{4b}{a}\left(\frac{x}{2}\sqrt{a^2 - x^2} + \frac{a^2}{2}\arcsin\frac{x}{a}\right)\Bigg|_0^a = \pi ab.
\]
我们也可以用参数方程
\[
\begin{cases}
x = a\cos t, \\
y = b\sin t,
\end{cases}
\quad 0 \leqslant t \leqslant 2\pi,
\]
用类似于例~2 的方法来求
\[
A = -\int_0^{2\pi} b\sin t\,\mathrm{d}(a\cos t) = ab\int_0^{2\pi} \sin^2 t\,\mathrm{d}t = \pi ab.
\]
当 $a = b$ 时，就是圆的面积 $\pi a^2$．
\end{proof}

从例~2 和例~3 看出，假设曲线是一条简单闭曲线 $\Gamma$（即除两个端点相连接外，其他地方不自相交），它由参数方程
\[
\begin{cases}
x = \varphi(t), \\
y = \psi(t),
\end{cases}
\quad \alpha \leqslant t \leqslant \beta
\]
给出．曲线在端点连接是指
\[
\varphi(\alpha) = \varphi(\beta), \qquad \psi(\alpha) = \psi(\beta).
\]

\begin{figure}[htbp]
\centering
\begin{minipage}[t]{0.45\textwidth}
\centering
\includegraphics[width=0.9\textwidth]{fig8-6.png}
\figcap{图 8-6}\label{fig:8-6}
\end{minipage}
\hfill
\begin{minipage}[t]{0.45\textwidth}
\centering
\includegraphics[width=0.9\textwidth]{fig8-7.png}
\figcap{图 8-7}\label{fig:8-7}
\end{minipage}
\end{figure}

如果 $\varphi, \psi$ 都是逐段有连续微商的，且与垂直于 $x$ 轴的直线相交至多两点，那么闭曲线 $\Gamma$ 所围起来的有界区域的面积为
\begin{equation}\label{eq:8-2-1}
S = -\int_\alpha^\beta \psi(t)\varphi'(t)\,\mathrm{d}t.
\end{equation}
事实上，这时曲线可由 $x = a, x = b$ 两条直线分成上下两段（见图~\ref{fig:8-7}），其方程分别可以写成
\[
\begin{aligned}
y &= y_2(x), \quad a \leqslant x \leqslant b, \\
y &= y_1(x), \quad a \leqslant x \leqslant b
\end{aligned}
\]
的形式，则面积
\begin{equation}\label{eq:8-2-2}
S = \int_a^b [y_2(x) - y_1(x)]\,\mathrm{d}x.
\end{equation}
设 $t = \alpha$ 对应于曲线的 $P$ 点（计算结果将表明，公式与 $P$ 点所处的位置无关），曲线与 $x = a$ 相交的点的对应参数为 $t_1$，与 $x = b$ 相交的点的对应参数为 $t_2$，则在面积公式\eqref{eq:8-2-2}中用换元积分法，令 $x = \varphi(t)$ 得到
\begin{align*}
S &= \int_a^b [y_2(x) - y_1(x)]\,\mathrm{d}x \\
&= \int_{t_1}^{t_2} \psi(t)\varphi'(t)\,\mathrm{d}t - \int_{t_1}^\beta \psi(t)\varphi'(t)\,\mathrm{d}t - \int_\alpha^{t_2} \psi(t)\varphi'(t)\,\mathrm{d}t \\
&= -\int_\alpha^{t_2} \psi(t)\varphi'(t)\,\mathrm{d}t - \int_{t_2}^{t_1} \psi(t)\varphi'(t)\,\mathrm{d}t - \int_{t_1}^\beta \psi(t)\varphi'(t)\,\mathrm{d}t \\
&= -\int_\alpha^\beta \psi(t)\varphi'(t)\,\mathrm{d}t.
\end{align*}
这便证明了\eqref{eq:8-2-1}式．记住，这时参数增加的方向为图上箭头所指的方向，这方向的特征是，当人沿曲线的方向行走时，曲线所围的有界区域（也就是我们要求面积的区域）永远在人的左方．我们就用这个法则规定曲线的正向．再说一遍，平面简单曲线的正向是指曲线的这样一个方向，当人沿着这个方向前进时，曲线所围的有界区域永远在人的左方．这样，我们可以通过 $y = \psi(t), \mathrm{d}x = \varphi'(t)\,\mathrm{d}t$，把\eqref{eq:8-2-1}式简化写成
\begin{equation}\label{eq:8-2-3}
S = -\int_\alpha^\beta \psi(t)\varphi'(t)\,\mathrm{d}t = -\int_\alpha^\beta \psi(t)\,\mathrm{d}\varphi(t) = -\oint_\Gamma y\,\mathrm{d}x.
\end{equation}
这公式的右边我们称为在曲线 $\Gamma$ 上的积分．其计算方法是用曲线的参数方程 $y = \psi(t), x = \varphi(t)$ 代进去，化为定积分，积分的下上限分别等于始点与终点所对应的参数值，参数由下限变到上限时正对应于曲线的正向（注意下限并不总是小于上限）．例如例~2 的曲线由参数方程
\[
\begin{cases}
x = a + a\cos t, \\
y = a + a\sin t,
\end{cases}
\quad 0 \leqslant t \leqslant 2\pi
\]
给出（见图~\ref{fig:8-5}）．这时 $t$ 由 $0$ 增加到 $2\pi$ 时，恰巧对应于曲线的正向．故
\begin{align*}
S &= -\oint y\,\mathrm{d}x \\
&= -\int_0^{2\pi} a(1 + \sin t)\,\mathrm{d}[a(1 + \cos t)] \\
&= \int_0^{2\pi} a^2(1 + \sin t)\sin t\,\mathrm{d}t,
\end{align*}
这正是我们在例~2 计算中推导过的，例~3 的情形完全类似（见图~\ref{fig:8-6}）．这时，曲线的参数方程为
\[
\begin{cases}
x = a\cos t, \\
y = b\sin t,
\end{cases}
\quad 0 \leqslant t \leqslant 2\pi.
\]
参数 $t$ 由 $0$ 变到 $2\pi$，正好也是曲线的正向．因此
\[
S = -\oint y\,\mathrm{d}x = -\int_0^{2\pi} b\sin t\,\mathrm{d}(a\cos t) = ab\int_0^{2\pi} \sin^2 t\,\mathrm{d}t.
\]
为了对称起见，我们把 $x$ 与 $y$ 的位置互换．假设简单闭曲线 $\Gamma$ 可以用两条直线 $y = c$ 与 $y = d$ 分成左右两段，其方程分别为
\[
x = g(y), \quad c \leqslant y \leqslant d,
\]
\[
x = h(y), \quad c \leqslant y \leqslant d.
\]

\begin{figure}[htbp]
\centering
\includegraphics[width=0.45\textwidth]{fig8-8.png}
\figcap{图 8-8}\label{fig:8-8}
\end{figure}

则曲线所围的有界区域的面积
\[
S = \int_c^d [g(y) - h(y)]\,\mathrm{d}y.
\]
假设曲线仍然用参数方程
\[
\begin{cases}
x = \varphi(t), \\
y = \psi(t),
\end{cases}
\quad \alpha \leqslant t \leqslant \beta
\]
给出，参数由 $\alpha$ 变到 $\beta$ 时正好表示曲线 $\Gamma$ 的正向，则用积分换元法则，类似于前面的推导得
\begin{equation}\label{eq:8-2-4}
S = \int_\alpha^\beta \varphi(t)\,\mathrm{d}\psi(t) = \oint x\,\mathrm{d}y.
\end{equation}
值得注意的是，和前面的公式相比，这里积分号前没有负号．当闭曲线所围面积同时可用\eqref{eq:8-2-3}，\eqref{eq:8-2-4}计算时，把两式相加除以 $2$，得
\[
S = \frac{1}{2}\oint (x\,\mathrm{d}y - y\,\mathrm{d}x).
\]
这就是平面简单闭曲线所围区域面积的计算公式，其计算方法是用曲线的参数方程 $x = \varphi(t), y = \psi(t)$ 代进去，下上限分别为曲线始点与终点所对应的参数．当参数由下限变到上限时，必须与曲线的正向相同，这公式在本教程下册中将要用到．


\subsection{极坐标下平面图形的面积}\label{subsec:8-2-2}

要求在极坐标下，由曲线 $r = r(\theta)$ 与向径 $\theta = \alpha$，$\theta = \beta$ 所围成的平面图形的面积，通常假设 $r = r(\theta)$ 连续，非负，$\beta - \alpha \leqslant 2\pi$．

我们用积分的定义来求．如图~\ref{fig:8-9} 所示，给 $[\alpha,\beta]$ 以分法
\[
\alpha = \theta_0 < \theta_1 < \cdots < \theta_n = \beta,
\]
看中间的一小块．设 $r(\theta)$ 在 $[\theta_{i-1},\theta_i]$ 的最大值、最小值分别为 $M_i, m_i$，则小块图形就夹在分别以 $M_i, m_i$ 为半径的两块圆扇形面积之间，因此
\[
\frac{1}{2}m_i^2\Delta\theta_i \leqslant \Delta A_i \leqslant \frac{1}{2}M_i^2\Delta\theta_i.
\]
对 $i$ 求和，由 $A = \sum_{i=1}^n \Delta A_i$，知
\[
\frac{1}{2}\sum_{i=1}^n m_i^2\Delta\theta_i \leqslant A \leqslant \frac{1}{2}\sum_{i=1}^n M_i^2\Delta\theta_i.
\]

\begin{figure}[htbp]
\centering
\begin{minipage}[t]{0.45\textwidth}
\centering
\includegraphics[width=0.9\textwidth]{fig8-9.png}
\figcap{图 8-9}\label{fig:8-9}
\end{minipage}
\hfill
\begin{minipage}[t]{0.45\textwidth}
\centering
\includegraphics[width=0.9\textwidth]{fig8-10.png}
\figcap{图 8-10}\label{fig:8-10}
\end{minipage}
\end{figure}

不等式左右两边是函数 $\frac{1}{2}r^2(\theta)$ 的黎曼和，令 $\lambda = \max_{1 \leqslant i \leqslant n} |\Delta\theta_i| \to 0$ 取极限，它们有相同的极限，故
\[
A = \frac{1}{2}\int_\alpha^\beta r^2(\theta)\,\mathrm{d}\theta.
\]
从而用微元的观点看
\[
\mathrm{d}A = \frac{1}{2}r^2(\theta)\,\mathrm{d}\theta,
\]
几何意义是很清楚的．

\begin{example}
求心脏线
\[
r = a(1 + \cos\theta), \quad 0 \leqslant \theta \leqslant 2\pi
\]
所围的面积，其中 $a > 0$（图~\ref{fig:8-10}）．
\end{example}

\begin{proof}[解]
\begin{align*}
A &= \frac{1}{2}\int_0^{2\pi} a^2(1 + \cos\theta)^2\,\mathrm{d}\theta \\
&= a^2\int_0^\pi (1 + 2\cos\theta + \cos^2\theta)\,\mathrm{d}\theta \\
&= a^2\pi + a^2\int_0^\pi \frac{1 + \cos 2\theta}{2}\,\mathrm{d}\theta \\
&= a^2\pi + \frac{a^2}{2}\pi = \frac{3}{2}\pi a^2.
\end{align*}
\end{proof}


\subsection{旋转体的体积}\label{subsec:8-2-3}

设有连续曲线 $y = f(x)$，满足
\[
f(x) \geqslant 0, \quad x \in [a,b].
\]
将此曲线绕 $x$ 轴旋转一周，求所产生的旋转体的体积 $V$．

给 $[a,b]$ 以分法
\[
a = x_0 < x_1 < \cdots < x_n = b,
\]
对应地便把旋转体分成 $n$ 片小块．记 $f(x)$ 在 $[x_{i-1}, x_i]$ 的最大值、最小值分别为 $M_i, m_i$，则第 $i$ 片旋转体夹在分别以 $M_i, m_i$ 为半径的两个圆柱体之间，因此它的体积 $\Delta V_i$ 满足
\[
\pi m_i^2\Delta x_i \leqslant \Delta V_i \leqslant \pi M_i^2\Delta x_i,
\]
求和得
\[
\pi\sum_{i=1}^n m_i^2\Delta x_i \leqslant V \leqslant \pi\sum_{i=1}^n M_i^2\Delta x_i.
\]
左、右两边都是函数 $\pi[f(x)]^2$ 的黎曼和，令 $\lambda = \max_{1 \leqslant i \leqslant n}\{\Delta x_i\} \to 0$，则左右两边的极限相等，故
\[
V = \pi\int_a^b [f(x)]^2\,\mathrm{d}x.
\]
从体积微元的观点看，
\[
\mathrm{d}V = \pi[f(x)]^2\,\mathrm{d}x = \pi y^2\,\mathrm{d}x,
\]
在局部看，旋转体的体积近似于圆柱体的体积，这公式是显然的（图~\ref{fig:8-11}）．

\begin{figure}[htbp]
\centering
\begin{minipage}[t]{0.45\textwidth}
\centering
\includegraphics[width=0.9\textwidth]{fig8-11.png}
\figcap{图 8-11}\label{fig:8-11}
\end{minipage}
\hfill
\begin{minipage}[t]{0.45\textwidth}
\centering
\includegraphics[width=0.9\textwidth]{fig8-12.png}
\figcap{图 8-12}\label{fig:8-12}
\end{minipage}
\end{figure}

\begin{example}
求星形线 $x^{2/3} + y^{2/3} = a^{2/3}$ ($a > 0$) 绕 $x$ 轴旋转所成之旋转体的体积（见图~\ref{fig:8-12}）．
\end{example}

\begin{proof}[解]
解出显式函数
\[
y^2 = (a^{2/3} - x^{2/3})^3,
\]
则旋转体体积为
\[
V = \pi\int_{-a}^a y^2\,\mathrm{d}x = \pi\int_{-a}^a (a^{2/3} - x^{2/3})^3\,\mathrm{d}x = \frac{32}{105}\pi a^3.
\]
\end{proof}

如果立体图形不是旋转体，但假设对每个 $x \in [a, b]$，我们都知道立体与过 $x$ 而垂直于 $x$ 轴的平面所截出的图形面积为 $A(x)$（见图~\ref{fig:8-13}），则我们很容易看出，立体的体积微元为
\[
\mathrm{d}V = A(x)\,\mathrm{d}x,
\]
从而体积为
\[
V = \int_a^b A(x)\,\mathrm{d}x.
\]

\begin{figure}[htbp]
\centering
\begin{minipage}[t]{0.45\textwidth}
\centering
\includegraphics[width=0.95\textwidth]{fig8-13.png}
\figcap{图 8-13}\label{fig:8-13}
\end{minipage}
\hfill
\begin{minipage}[t]{0.45\textwidth}
\centering
\includegraphics[width=0.8\textwidth]{fig8-14.png}
\figcap{图 8-14}\label{fig:8-14}
\end{minipage}
\end{figure}

\begin{example}
求半径为 $a$ 的圆柱，被一与底面成 $\alpha$ 角而过底面直径的平面所截出的楔形体的体积．
\end{example}

\begin{proof}[解]
如图~\ref{fig:8-14} 所示，取底面为 $Oxy$ 坐标平面，斜面与底面的交线为 $x$ 轴，圆柱与水平面交线为圆周，其方程为
\[
y = \sqrt{a^2 - x^2}, \quad -a \leqslant x \leqslant a.
\]
过点 $x$ 且垂直于 $x$ 轴的平面与楔形体交一三角形，其底角为 $\alpha$，底边长为 $\sqrt{a^2 - x^2}$，故三角形的面积
\[
A(x) = \frac{1}{2}(a^2 - x^2)\tan\alpha;
\]
因此
\[
V = \int_{-a}^a A(x)\,\mathrm{d}x = \frac{1}{2}\tan\alpha\int_{-a}^a (a^2 - x^2)\,\mathrm{d}x = \tan\alpha\left(a^2 x - \frac{x^3}{3}\right)\Bigg|_0^a = \frac{2}{3}a^3\tan\alpha.
\]
\end{proof}


\subsection{曲线的弧长}\label{subsec:8-2-4}

由方程
\[
\begin{cases}
x = \varphi(t), \\
y = \psi(t),
\end{cases}
\quad \alpha \leqslant t \leqslant \beta
\]
决定的 $(x,y)$ 在坐标平面上构成的点集 $\Gamma$，称为曲线，如果 $\varphi, \psi$ 是连续函数，且点集 $\Gamma$ 没有重点，即对任意 $t_1 \neq t_2$，有
\[
\big(\varphi(t_1), \psi(t_1)\big) \neq \big(\varphi(t_2), \psi(t_2)\big).
\]
但端点可以例外，即可以有 $\varphi(\alpha) = \varphi(\beta), \psi(\alpha) = \psi(\beta)$，这时 $\Gamma$ 称为闭曲线．

注意，曲线作为平面的特殊点集，它是有方向的．$t$ 从 $\alpha$ 变到 $\beta$，决定曲线的一个方向；$t$ 从 $\beta$ 变到 $\alpha$，对应于曲线的另一方向．

$\Gamma$ 作为一条曲线，它的长度怎样求呢？类似于圆周长用圆内接正多边形的周长当边数趋向于无穷时的极限来定义，曲线的长度也可定义为内接折线长的极限．

给 $[\alpha,\beta]$ 任意一个分法：
\[
\alpha = t_0 < t_1 < \cdots < t_n = \beta,
\]
这时对应于曲线上从 $A$ 到 $B$ 有顺序的 $n+1$ 个点（图~\ref{fig:8-15}）
\[
A = M_0, M_1, M_2, \cdots, M_n = B,
\]
其中
\[
M_k = \big(\varphi(t_k), \psi(t_k)\big).
\]
记
\[
p = \sum_{k=1}^n M_{k-1}M_k,
\]
其中 $M_{k-1}M_k$ 表示直线段 $M_{k-1}M_k$ 的长度．记 $\lambda = \max_{1 \leqslant k \leqslant n} \{\Delta t_k\}$，其中 $\Delta t_k = t_k - t_{k-1}$，若当 $\lambda \to 0$ 时，$p$ 的极限存在，即
\[
\lim_{\lambda\to 0} p = s,
\]
则称曲线 $\Gamma$ 是可求长的，而 $s$ 称为 $\Gamma$ 的弧长．

\begin{figure}[htbp]
\centering
\includegraphics[width=0.45\textwidth]{fig8-15.png}
\figcap{图 8-15}\label{fig:8-15}
\end{figure}

极限的正确含义自然是：对任意 $\varepsilon > 0$，存在 $\delta > 0$，不管分法如何，只要 $\lambda < \delta$，则相应的折线长度 $p$ 满足
\[
|p - s| < \varepsilon.
\]
如果 $\varphi', \psi'$ 在 $[\alpha,\beta]$ 连续，且
\[
[\varphi'(t)]^2 + [\psi'(t)]^2 \neq 0, \quad \forall\; t \in [\alpha,\beta],
\]
则称曲线 $\Gamma$ 是光滑的．下面我们来证明光滑曲线是可求长的．

记 $M_k = (x_k, y_k)$，其中 $x_k = \varphi(t_k), y_k = \psi(t_k)$，则由微分中值定理知，存在 $\xi_k \in [t_{k-1}, t_k], \eta_k \in [t_{k-1}, t_k]$，使得
\[
\Delta x_k = x_k - x_{k-1} = \varphi(t_k) - \varphi(t_{k-1}) = \varphi'(\xi_k)\Delta t_k,
\]
\[
\Delta y_k = y_k - y_{k-1} = \psi(t_k) - \psi(t_{k-1}) = \psi'(\eta_k)\Delta t_k.
\]
这样
\[
M_{k-1}M_k = \sqrt{\Delta x_k^2 + \Delta y_k^2} = \sqrt{[\varphi'(\xi_k)]^2 + [\psi'(\eta_k)]^2}\,\Delta t_k,
\]
从而
\[
p = \sum_{k=1}^n M_{k-1}M_k = \sum_{k=1}^n \sqrt{[\varphi'(\xi_k)]^2 + [\psi'(\eta_k)]^2}\,\Delta t_k.
\]
由于 $\varphi', \psi'$ 连续，根据连续函数的可积性知，当 $\lambda \to 0$ 时黎曼和
\[
\sigma = \sum_{k=1}^n \sqrt{[\varphi'(\xi_k)]^2 + [\psi'(\xi_k)]^2}\,\Delta t_k
\]
的极限存在，即 $\lim_{\lambda\to 0}\sigma$ 存在，且等于
\[
\int_\alpha^\beta \sqrt{[\varphi'(t)]^2 + [\psi'(t)]^2}\,\mathrm{d}t.
\]
如果证得
\[
\lim_{\lambda\to 0}(p - \sigma) = 0,
\]
则当 $\lambda \to 0$ 时折线长 $p$ 的极限存在．为此，记 $M_k, m_k$ 分别为 $\psi'(t)$ 在 $[t_{k-1}, t_k]$ 的最大值与最小值，则由一致连续性知，任给 $\varepsilon > 0$，存在 $\delta > 0$，只要 $\lambda < \delta$，就有
\[
M_k - m_k < \frac{\varepsilon}{\beta - \alpha}, \qquad \forall\; k = 1, 2, \cdots, n.
\]
这样，只要 $\lambda < \delta$，就有
\begin{align*}
|p - \sigma| &\leqslant \sum_{k=1}^n \left|\sqrt{[\varphi'(\xi_k)]^2 + [\psi'(\eta_k)]^2} - \sqrt{[\varphi'(\xi_k)]^2 + [\psi'(\xi_k)]^2}\right|\Delta t_k \\
&\leqslant \sum_{k=1}^n |\psi'(\eta_k) - \psi'(\xi_k)|\Delta t_k \leqslant \sum_{k=1}^n (M_k - m_k)\Delta t_k \\
&< \frac{\varepsilon}{\beta - \alpha}\sum_{k=1}^n \Delta t_k = \varepsilon,
\end{align*}
其中在不等式放大中用到了下面的事实：
\begin{align*}
\left|\sqrt{A^2 + B^2} - \sqrt{A^2 + B_1^2}\right| &= \frac{|B^2 - B_1^2|}{\sqrt{A^2 + B^2} + \sqrt{A^2 + B_1^2}} \\
&\leqslant \frac{|B + B_1|}{\sqrt{B^2} + \sqrt{B_1^2}}|B - B_1| = \frac{|B + B_1|}{|B| + |B_1|}|B - B_1| \leqslant |B - B_1|.
\end{align*}
这样，我们就证明了
\[
\lim_{\lambda\to 0}(p - \sigma) = 0,
\]
从而极限
\[
\lim_{\lambda\to 0}p = \lim_{\lambda\to 0}(p - \sigma + \sigma) = \lim_{\lambda\to 0}(p - \sigma) + \lim_{\lambda\to 0}\sigma = \lim_{\lambda\to 0}\sigma,
\]
这就是说 $\Gamma$ 是可求长的，且长度
\[
s = \int_\alpha^\beta \sqrt{[\varphi'(t)]^2 + [\psi'(t)]^2}\,\mathrm{d}t.
\]
如果曲线由 $y = f(x)$ ($a \leqslant x \leqslant b$) 给出，$f'(x)$ 在 $[a, b]$ 连续，则曲线是可求长的．因为把它写成参数方程
\[
\begin{cases}
y = f(x), \quad a \leqslant x \leqslant b, \\
x = x,
\end{cases}
\]
根据前面的公式便得曲线的弧长
\[
s = \int_a^b \sqrt{1 + [f'(x)]^2}\,\mathrm{d}x.
\]

取定曲线上的一点，例如 $A = M_0$，对于曲线上任意一点 $M = (\varphi(t), \psi(t))$，它对应于参数 $t \in [\alpha, \beta]$，则曲线 $\wideparen{M_0M}$ 的弧长便是 $t$ 的函数，记为 $s(t)$，它等于
\[
s(t) = \int_\alpha^t \sqrt{[\varphi'(\tau)]^2 + [\psi'(\tau)]^2}\,\mathrm{d}\tau.
\]
对 $t$ 求微分，便得
\[
\mathrm{d}s = \sqrt{[\varphi'(t)]^2 + [\psi'(t)]^2}\,\mathrm{d}t,
\]
称为弧微分．取平方，注意到 $\mathrm{d}x = \varphi'(t)\mathrm{d}t, \mathrm{d}y = \psi'(t)\mathrm{d}t$，便有
\[
\mathrm{d}s^2 = \left([\varphi'(t)]^2 + [\psi'(t)]^2\right)\mathrm{d}t^2 = \mathrm{d}x^2 + \mathrm{d}y^2,
\]
它的几何解释如图~\ref{fig:8-16}，$\mathrm{d}s$ 就是切线的长度，它是 $\mathrm{d}t$ 的线性函数，与弧长相差一个较 $\mathrm{d}t$ 高阶的无穷小量．由 $\mathrm{d}x, \mathrm{d}y, \mathrm{d}s$ 构成的三角形称为微分三角形．上面的等式是微分三角形的勾股弦定理．

\begin{figure}[htbp]
\centering
\begin{minipage}[t]{0.45\textwidth}
\centering
\includegraphics[width=0.85\textwidth]{fig8-16.png}
\figcap{图 8-16}\label{fig:8-16}
\end{minipage}
\hfill
\begin{minipage}[t]{0.45\textwidth}
\centering
\includegraphics[width=0.95\textwidth]{fig8-17.png}
\figcap{图 8-17}\label{fig:8-17}
\end{minipage}
\end{figure}

在直角坐标系中，有
\[
\mathrm{d}s^2 = (1 + y'^2)\mathrm{d}x^2.
\]

\begin{example}
求旋轮线一拱的弧长（图~\ref{fig:8-17}）．在圆周上取一定点 $M$，当圆沿一直线无滑动地向前滚动时，$M$ 的轨迹称为旋轮线．
\end{example}

\begin{proof}[解]
取圆上定点 $M$ 与直线相切时的切点作原点，直线为 $x$ 轴．设圆周滚动了弧长 $at$，则 $OA = \wideparen{MA} = at$，故 $M$ 的坐标 $(x, y)$ 应满足
\begin{align*}
x &= at - a\sin t = a(t - \sin t), \\
y &= a - a\cos t = a(1 - \cos t).
\end{align*}
对应于一拱，参数为 $0 \leqslant t \leqslant 2\pi$，这时
\[
\begin{cases}
x'(t) = a(1 - \cos t), \\
y'(t) = a\sin t,
\end{cases}
\]
因此
\[
[x'(t)]^2 + [y'(t)]^2 = a^2(1 + \cos^2 t - 2\cos t + \sin^2 t) = 2a^2(1 - \cos t),
\]
故旋轮线一拱的长度为
\[
s = \int_0^{2\pi} a\sqrt{2(1 - \cos t)}\,\mathrm{d}t = 2a\int_0^{2\pi} \sin\frac{t}{2}\,\mathrm{d}t = 8a.
\]
\end{proof}

如果曲线由极坐标 $r = r(\theta)$ ($\alpha \leqslant \theta \leqslant \beta$) 给出，则相应的参数方程是
\[
\begin{cases}
x = r(\theta)\cos\theta, \\
y = r(\theta)\sin\theta,
\end{cases}
\quad \alpha \leqslant \theta \leqslant \beta.
\]
这时
\[
x'_\theta = r'(\theta)\cos\theta - r(\theta)\sin\theta,
\]
\[
y'_\theta = r'(\theta)\sin\theta + r(\theta)\cos\theta,
\]
因此
\[
(x'_\theta)^2 + (y'_\theta)^2 = [r'(\theta)]^2 + [r(\theta)]^2,
\]
故曲线弧长在极坐标下的公式应为
\[
s = \int_\alpha^\beta \sqrt{[r(\theta)]^2 + [r'(\theta)]^2}\,\mathrm{d}\theta.
\]
相应地，弧微分公式为
\begin{align*}
\mathrm{d}s^2 &= r^2\,\mathrm{d}\theta^2 + r'^2\,\mathrm{d}\theta^2 \\
&= r^2\,\mathrm{d}\theta^2 + \mathrm{d}r^2,
\end{align*}
从几何上看，这也是“微分三角形”的勾股弦定理（图~\ref{fig:8-18}）．

\begin{figure}[htbp]
\centering
\begin{minipage}[t]{0.45\textwidth}
\centering
\includegraphics[width=0.85\textwidth]{fig8-18.png}
\figcap{图 8-18}\label{fig:8-18}
\end{minipage}
\hfill
\begin{minipage}[t]{0.45\textwidth}
\centering
\includegraphics[width=0.95\textwidth]{fig8-19.png}
\figcap{图 8-19}\label{fig:8-19}
\end{minipage}
\end{figure}

\begin{example}
求双扭线
\[
r^2 = 2a^2\cos 2\theta \qquad (a > 0)
\]
从 $\theta = 0$ 到 $\theta = \frac{\pi}{6}$ 的弧长（图~\ref{fig:8-19}）．
\end{example}

\begin{proof}[解]
在方程 $r^2 = 2a^2\cos 2\theta$ 两端对 $\theta$ 求微商，有
\[
2r r' = -4a^2\sin 2\theta,
\]
因此
\[
r' = -\frac{2a^2\sin 2\theta}{r},
\]
故
\begin{align*}
\sqrt{[r(\theta)]^2 + [r'(\theta)]^2} &= \sqrt{r^2 + \frac{4a^4\sin^2 2\theta}{r^2}} \\
&= \sqrt{\frac{r^4 + 4a^4\sin^2 2\theta}{r^2}} = \frac{\sqrt{2}a}{\sqrt{\cos 2\theta}} = \frac{\sqrt{2}a}{\sqrt{1 - 2\sin^2\theta}},
\end{align*}
于是
\[
s = \int_0^{\pi/6}\frac{\sqrt{2}a}{\sqrt{1 - 2\sin^2\theta}}\,\mathrm{d}\theta = \sqrt{2}a\int_0^{\pi/6}\frac{\mathrm{d}\theta}{\sqrt{1 - 2\sin^2\theta}}.
\]
这个被积函数的原函数不是初等函数，我们把这种积分称为“椭圆积分”．
\end{proof}

我们很容易把上述讨论推广到空间曲线上．设在空间中给定参数方程
\[
\begin{cases}
x = \varphi(t), \\
y = \psi(t), \\
z = \chi(t),
\end{cases}
\quad \alpha \leqslant t \leqslant \beta,
\]
其中 $\varphi, \psi, \chi$ 是连续的，则全体值域在三维空间中构成的点集 $\Gamma$ 称为空间的一条曲线（图~\ref{fig:8-20}）．当然，我们也要求它是不自相交的，只有在端点 $\alpha, \beta$ 例外．类似地可以给出 $\Gamma$ 可求长的定义并证明，光滑曲线是可求长的，且其弧长
\[
s = \int_\alpha^\beta \sqrt{[\varphi'(t)]^2 + [\psi'(t)]^2 + [\chi'(t)]^2}\,\mathrm{d}t.
\]

\begin{figure}[htbp]
\centering
\begin{minipage}[t]{0.45\textwidth}
\centering
\includegraphics[width=0.9\textwidth]{fig8-20.png}
\figcap{图 8-20}\label{fig:8-20}
\end{minipage}
\hfill
\begin{minipage}[t]{0.45\textwidth}
\centering
\includegraphics[width=0.9\textwidth]{fig8-21.png}
\figcap{图 8-21}\label{fig:8-21}
\end{minipage}
\end{figure}

类似地有弧微分，还可以得到
\begin{align*}
\mathrm{d}s^2 &= \varphi'^2\mathrm{d}t^2 + \psi'^2\mathrm{d}t^2 + \chi'^2\mathrm{d}t^2 \\
&= \mathrm{d}x^2 + \mathrm{d}y^2 + \mathrm{d}z^2.
\end{align*}

曲线也可以用向量来表示．空间的向量记为 $\boldsymbol{r}$，它有三个分量 $\boldsymbol{r} = (x, y, z)$．如果对于任意 $t$，有向量 $\boldsymbol{r}(t)$ 与之对应，这就是向量值的函数．$\boldsymbol{r}(t)$ 给定当且仅当它的三个分量 $x(t), y(t), z(t)$ 给定．我们说 $\boldsymbol{r}(t)$ 在 $t_0$ 是连续的，如果
\[
\lim_{t\to t_0}\boldsymbol{r}(t) = \boldsymbol{r}(t_0).
\]
这意思显然是指，对任给 $\varepsilon > 0$，存在 $\delta > 0$，只要 $|t - t_0| < \delta$，便有
\[
|\boldsymbol{r}(t) - \boldsymbol{r}(t_0)| < \varepsilon,
\]
其中 $\boldsymbol{r}(t) - \boldsymbol{r}(t_0)$ 是向量的减法，$|\boldsymbol{r}|$ 表示向量 $\boldsymbol{r}$ 的长度．显然的结论是，$\boldsymbol{r}(t)$ 在 $t_0$ 连续，当且仅当它的三个分量 $x(t), y(t), z(t)$ 在 $t_0$ 连续．当 $\boldsymbol{r}(t)$ 连续时，
\[
\boldsymbol{r} = \boldsymbol{r}(t)
\]
自然就给出空间曲线的方程（图~\ref{fig:8-21}）．如果把 $t$ 看作时间，这就是质点的运动方程．

对向量 $\boldsymbol{r}(t)$，自然也可以定义它的微商：
\[
\frac{\mathrm{d}\boldsymbol{r}(t)}{\mathrm{d}t} = \lim_{\Delta t\to 0}\frac{\boldsymbol{r}(t+\Delta t) - \boldsymbol{r}(t)}{\Delta t}.
\]
按定义，$\boldsymbol{r}(t)$ 的微商也是向量，它的三个分量分别就是原来的三个分量的微商：
\[
\frac{\mathrm{d}\boldsymbol{r}(t)}{\mathrm{d}t} = \left(\frac{\mathrm{d}x(t)}{\mathrm{d}t}, \frac{\mathrm{d}y(t)}{\mathrm{d}t}, \frac{\mathrm{d}z(t)}{\mathrm{d}t}\right) = \big(x'(t), y'(t), z'(t)\big).
\]
如果把 $\boldsymbol{r} = \boldsymbol{r}(t)$ 看成质点的运动方程，由定义自然就知道其微商 $\frac{\mathrm{d}\boldsymbol{r}(t)}{\mathrm{d}t}$ 就是质点的运动速度（向量）$\boldsymbol{v}(t)$，二阶微商 $\frac{\mathrm{d}^2\boldsymbol{r}(t)}{\mathrm{d}t^2}$ 就是加速度（向量）$\boldsymbol{a}(t)$，因此
\[
|\boldsymbol{v}(t)| = \sqrt{[x'(t)]^2 + [y'(t)]^2 + [z'(t)]^2},
\]
故弧长
\[
s = \int_\alpha^\beta \sqrt{[x'(t)]^2 + [y'(t)]^2 + [z'(t)]^2}\,\mathrm{d}t = \int_\alpha^\beta |\boldsymbol{v}(t)|\,\mathrm{d}t.
\]
这说明，弧长正好等于质点在时间间隔 $[\alpha,\beta]$ 中所走过的路程（长度），这和人们的直观完全符合．

本段对曲线的讨论，都是通过曲线的参数方程给出的．参数的选择可以多种多样，但有一种参数的选取却是由曲线本身给出的，这就是弧长．事实上，设曲线是光滑曲线，则
\[
s(t) = \int_\alpha^t \sqrt{[x'(\tau)]^2 + [y'(\tau)]^2 + [z'(\tau)]^2}\,\mathrm{d}\tau,
\]
这时
\[
s'(t) = \sqrt{[x'(t)]^2 + [y'(t)]^2 + [z'(t)]^2} > 0.
\]
这说明 $s(t)$ 是严格增函数，因此有反函数 $t = t(s)$ 存在，把它代回原来曲线的参数方程
\[
\begin{cases}
x = \varphi(t) = \varphi\big(t(s)\big), \\
y = \psi(t) = \psi\big(t(s)\big), \\
z = \chi(t) = \chi\big(t(s)\big),
\end{cases}
\quad 0 \leqslant s \leqslant l,
\]
其中 $l$ 是总弧长．这样，曲线就通过弧长作参数给出来，我们把曲线用弧长作参数给出的参数方程，称为曲线的本性方程．上面的过程给出了由其他参数方程求本性方程的过程，其逆步骤也就给出了由本性方程求其他参数方程的做法．


\subsection{曲线的曲率}\label{subsec:8-2-5}

我们只讲平面曲线的情形．

曲线的曲率描述的是曲线弯曲的程度．许多实际问题需要考虑曲线的弯曲程度，例如在工程技术中，往往遇到梁或轴因受外力作用而弯曲变形，为了保证使用安全，在设计时，必须对弯曲程度有所了解，以便将它限制在一定的范围之内．又如铁道设计，在拐弯的时候就不能让其弯曲程度太大，否则，火车在运行时就会出问题，等等．

究竟如何来刻画曲线的弯曲程度呢？考察两条长度相同的曲线 $\wideparen{AB}$ 和 $\wideparen{A'B}$（图~\ref{fig:8-22}）．当动点从点 $A$ 沿 $\wideparen{AB}$ 运动到点 $B$ 时，切线 $\tau_A$ 也随着转动到切线 $\tau_B$．记 $\Delta\varphi$ 为这两条切线之间的夹角．从图形看出，它们正巧等于切线 $\tau_B$ 和 $x$ 轴的交角与切线 $\tau_A$ 和 $x$ 轴的交角之差．同样，另一条曲线 $\wideparen{A'B}$ 的两个端点的切线 $\tau'_A, \tau'_B$ 的夹角为 $\Delta\varphi'$．显然，切线间夹角愈大，曲线弯曲的程度愈大．这里，$\Delta\varphi' > \Delta\varphi$，故 $\wideparen{A'B}$ 的弯曲程度比 $\wideparen{AB}$ 的弯曲程度大．

\begin{figure}[htbp]
\centering
\includegraphics[width=0.65\textwidth]{fig8-22.png}
\figcap{图 8-22}\label{fig:8-22}
\end{figure}

不过，切线的夹角大小还不能完全刻画曲线段的弯曲程度，还得看这两条切线间弧的长度有多大．对同一夹角，弧长愈短的，弯曲程度就愈大，如图~\ref{fig:8-23} 所表示的，切线夹角 $\Delta\varphi$ 相同，但 $\wideparen{A'B'}$ 的弧长要比 $\wideparen{AB}$ 的弧长短，因此 $\wideparen{A'B'}$ 的弯曲程度要比 $\wideparen{AB}$ 的大．

由此可见，曲线段的弯曲程度与两个端点处切线的夹角成正比，与曲线段的长度成反比．故如果曲线段 $\wideparen{AB}$ 的长度为 $\Delta s$，端点 $A, B$ 处切线的夹角为 $\Delta\varphi$，那么称
\[
\bar{K} = \left|\frac{\Delta\varphi}{\Delta s}\right|
\]
为曲线段 $\wideparen{AB}$ 的平均曲率．它刻画了曲线段 $\wideparen{AB}$ 的平均弯曲程度．

对于半径为 $R$ 的圆来说（图~\ref{fig:8-24}），圆周上任意一段弧 $\wideparen{AB}$ 的切线间交角即切线方向变化的角度 $\Delta\varphi$ 等于 $OA$ 与 $OB$ 之间的夹角即圆心角 $\Delta\alpha$．而 $\wideparen{AB}$ 的长度 $\Delta s = R\Delta\alpha$，因此弧 $\wideparen{AB}$ 的平均曲率为
\[
\bar{K} = \left|\frac{\Delta\varphi}{\Delta s}\right| = \frac{\Delta\alpha}{\Delta s} = \frac{1}{R}.
\]
这说明圆周上每段弧的平均曲率都一样，它们都等于 $\frac{1}{R}$．这与同一半径的圆，圆上的每点的弯曲程度都一样的直观是符合的．而对不同的圆，半径愈小，弯曲程度就愈大．

对于直线来说，因为它在每点的切线都一样，所以恒有 $\Delta\varphi = 0$，因此
\[
\bar{K} = \left|\frac{\Delta\varphi}{\Delta s}\right| = 0.
\]
这表示直线上任意一段的平均曲率都是 $0$，也就是说“直线不曲”，弯曲程度为 $0$．

\begin{figure}[htbp]
\centering
\begin{minipage}[t]{0.45\textwidth}
\centering
\includegraphics[width=0.7\textwidth]{fig8-23.png}
\figcap{图 8-23}\label{fig:8-23}
\end{minipage}
\hfill
\begin{minipage}[t]{0.45\textwidth}
\centering
\includegraphics[width=0.75\textwidth]{fig8-24.png}
\figcap{图 8-24}\label{fig:8-24}
\end{minipage}
\end{figure}

对于非圆非直线的曲线来说，每个地方的弯曲程度可能是不一样的，平均曲率只描写了曲线在这一段的“平均弯曲程度”．类似于引入瞬时速度概念的情形，$B$ 愈接近于 $A$，即 $\Delta s$ 愈小，$\wideparen{AB}$ 弧的平均曲率就能愈精确刻画曲线在 $A$ 处的弯曲程度，因此自然就定义当 $\Delta s \to 0$ 时平均曲率的极限为曲线在 $A$ 点的曲率：
\[
K = \left|\frac{\mathrm{d}\varphi}{\mathrm{d}s}\right| = \lim_{\Delta s\to 0}\left|\frac{\Delta\varphi}{\Delta s}\right|.
\]
曲率 $K$ 刻画了曲线在一点的弯曲程度．这里取绝对值是为了使曲率不变成负数．另外，我们也称曲率的倒数 $\frac{1}{K}$ 为曲线在该点的曲率半径．

下面我们推导计算曲率的公式．设曲线由参数方程
\[
\begin{cases}
x = x(t), \\
y = y(t),
\end{cases}
\quad \alpha \leqslant t \leqslant \beta
\]
给出，它们有二阶微商，注意对每个 $t \in [\alpha, \beta]$，曲线在对应点的切线斜率为 $\frac{\mathrm{d}y}{\mathrm{d}x}$，即切线与 $x$ 轴交角 $\varphi$ 的正切：
\[
\tan\varphi = \frac{\mathrm{d}y}{\mathrm{d}x} = \frac{y'(t)}{x'(t)},
\]
或
\[
\varphi = \arctan\frac{y'}{x'}.
\]
两边对 $t$ 求微商
\[
\frac{\mathrm{d}\varphi}{\mathrm{d}t} = \frac{1}{1 + \left(\frac{y'}{x'}\right)^2}\cdot\frac{y''x' - y'x''}{x'^2} = \frac{y''x' - y'x''}{x'^2 + y'^2}.
\]
而由弧微分知
\[
\frac{\mathrm{d}s}{\mathrm{d}t} = \sqrt{x'^2 + y'^2},
\]
因此
\[
K = \left|\frac{\mathrm{d}\varphi}{\mathrm{d}s}\right| = \frac{\left|\frac{\mathrm{d}\varphi}{\mathrm{d}t}\right|}{\left|\frac{\mathrm{d}s}{\mathrm{d}t}\right|} = \frac{|y''x' - y'x''|}{(x'^2 + y'^2)^{3/2}}.
\]
如果曲线由直角坐标 $y = f(x)$ 表出，那么这时 $x' = 1, x'' = 0$，故
\[
K = \frac{|y''|}{(1 + y'^2)^{3/2}}.
\]

\begin{example}
求半径为 $R$ 的圆周的曲率．
\end{example}

\begin{proof}[解]
\[
\begin{cases}
x = R\cos t, \\
y = R\sin t.
\end{cases}
\]
这时
\[
x' = -R\sin t, \quad y' = R\cos t,
\]
\[
x'' = -R\cos t, \quad y'' = -R\sin t.
\]
因此
\[
x'^2 + y'^2 = R^2, \quad y''x' - y'x'' = R^2,
\]
故
\[
K = \frac{R^2}{R^3} = \frac{1}{R}.
\]
即圆周的曲率处处相等，不同半径的圆周，曲率与半径成反比．
\end{proof}

\begin{example}
求悬链线
\[
y = \frac{a}{2}\left(\mathrm{e}^{x/a} + \mathrm{e}^{-x/a}\right)
\]
的曲率，其中 $a > 0$．
\end{example}

\begin{proof}[解]
\[
\sqrt{1 + y'^2} = \sqrt{1 + \frac{1}{4}\left(\mathrm{e}^{x/a} - \mathrm{e}^{-x/a}\right)^2} = \frac{y}{a},
\]
而易知
\[
y'' = \frac{y}{a^2},
\]
故
\[
K = \frac{y''}{(1 + y'^2)^{3/2}} = \frac{a}{y^2}.
\]
\end{proof}


\subsection{旋转体的侧面积}\label{subsec:8-2-6}

到现在为止，我们只会求一些可展成平面的曲面的面积，如圆柱面，圆锥面，圆台面等等．求一般曲面的面积是个很复杂的问题，我们留待以后再讲．现在讲一类特殊曲面的面积求法，它们就是旋转曲面．

设 $\wideparen{AB}$ 是光滑的平面曲线，绕 $x$ 轴旋转得一旋转曲面，我们要求它的侧面积（即除上下底以外的面积）．

\begin{figure}[htbp]
\centering
\begin{minipage}[t]{0.45\textwidth}
\centering
\includegraphics[width=0.9\textwidth]{fig8-25.png}
\figcap{图 8-25}\label{fig:8-25}
\end{minipage}
\hfill
\begin{minipage}[t]{0.45\textwidth}
\centering
\includegraphics[width=0.85\textwidth]{fig8-26.png}
\figcap{图 8-26}\label{fig:8-26}
\end{minipage}
\end{figure}

设 $\wideparen{AB}$ 光滑，参数 $t \in [\alpha, \beta]$，给 $[\alpha,\beta]$ 一个分法：$\alpha = t_0 < t_1 < \cdots < t_n = \beta$，对应地得到曲线 $\wideparen{AB}$ 上的点
\[
A = M_0, M_1, M_2, \cdots, M_n = B.
\]
用直线联结 $M_{k-1}M_k$，它绕 $x$ 轴旋转，得圆台面侧面积 $\Delta F_k$．若当 $\lambda = \max_{1 \leqslant i \leqslant n}|\Delta t_i| \to 0$ 时，极限
\[
\lim_{\lambda\to 0}\sum_{k=1}^n \Delta F_k = S,
\]
则称 $S$ 为旋转体的侧面积（图~\ref{fig:8-25}）．

下面我们来推导旋转体侧面积的公式．设曲线 $\wideparen{AB}$ 的本性方程为
\[
\begin{cases}
x = x(s), \\
y = y(s),
\end{cases}
\quad 0 \leqslant s \leqslant l,
\]
其中 $s$ 是 $\wideparen{AB}$ 的弧长参数，$l$ 是 $\wideparen{AB}$ 的弧长，$y(s) \geqslant 0$．给 $[0, l]$ 以分法：
\[
0 = s_0 < s_1 < \cdots < s_n = l,
\]
这时对应地得到曲线上的分点
\[
A = M_0, M_1, \cdots, M_n = B.
\]
$M_{k-1}M_k$ 绕 $x$ 轴旋转所得圆台面侧面积为
\[
\Delta F_k = 2\pi\frac{y_{k-1} + y_k}{2}\overline{M_{k-1}M_k},
\]
它们的和为
\[
F = \sum_{k=1}^n \Delta F_k = \pi\sum_{k=1}^n (y_{k-1} + y_k)\overline{M_{k-1}M_k}.
\]
记
\begin{align*}
\sigma &= \pi\sum_{k=1}^n (y_{k-1} + y_k)\Delta s_k \\
&= \pi\sum_{k=1}^n y(s_{k-1})\Delta s_k + \pi\sum_{k=1}^n y(s_k)\Delta s_k,
\end{align*}
其中 $\Delta s_k = s_k - s_{k-1}$ 为 $\wideparen{M_{k-1}M_k}$ 的弧长．由积分定义知，当 $\lambda = \max\{\Delta s_k\} \to 0$ 时，
\[
\lim_{\lambda\to 0}\sigma = 2\pi\int_0^l y(s)\,\mathrm{d}s.
\]
令 $M = \max_{0 \leqslant s \leqslant l}|y(s)|$，则
\begin{align*}
|F - \sigma| &= \pi\left|\sum_{k=1}^n (y_{k-1} + y_k)(\overline{M_{k-1}M_k} - \Delta s_k)\right| \\
&\leqslant \pi\sum_{k=1}^n (|y_{k-1}| + |y_k|)(\Delta s_k - \overline{M_{k-1}M_k}) \\
&\leqslant 2\pi M\left(l - \sum_{k=1}^n \overline{M_{k-1}M_k}\right) \to 0 \quad (\lambda \to 0).
\end{align*}
因此旋转体侧面积
\[
S = \lim_{\lambda\to 0}F = \lim_{\lambda\to 0}(F - \sigma) + \lim_{\lambda\to 0}\sigma = \lim_{\lambda\to 0}\sigma = \int_0^l 2\pi y(s)\,\mathrm{d}s.
\]
若曲线由
\[
\begin{cases}
x = \varphi(t), \\
y = \psi(t),
\end{cases}
\quad \alpha \leqslant t \leqslant \beta
\]
给出，这时弧长
\[
s = \int_\alpha^t \sqrt{[\varphi'(\tau)]^2 + [\psi'(\tau)]^2}\,\mathrm{d}\tau,
\]
\[
\mathrm{d}s = \sqrt{[\varphi'(t)]^2 + [\psi'(t)]^2}\,\mathrm{d}t,
\]
代入上式，得
\[
S = \int_\alpha^\beta 2\pi\psi(t)\sqrt{[\varphi'(t)]^2 + [\psi'(t)]^2}\,\mathrm{d}t.
\]
若曲线由直角坐标
\[
y = f(x), \quad a \leqslant x \leqslant b
\]
给出，则
\[
\mathrm{d}s = \sqrt{1 + [f'(x)]^2}\,\mathrm{d}x,
\]
因此旋转体侧面积为
\[
S = 2\pi\int_a^b f(x)\sqrt{1 + [f'(x)]^2}\,\mathrm{d}x.
\]

\begin{example}
求心脏线 $r = a(1 + \cos\theta)$ ($a > 0$) 绕 $x$ 轴旋转所成旋转体的侧面积（图~\ref{fig:8-26}）．
\end{example}

\begin{proof}[解]
$r = a(1 + \cos\theta)$ 是极坐标方程，可以认为它的参数方程是
\[
\begin{cases}
x = r(\theta)\cos\theta = a(1 + \cos\theta)\cos\theta, \\
y = r(\theta)\sin\theta = a(1 + \cos\theta)\sin\theta.
\end{cases}
\]
旋转体可以认为是由心脏线的上半支绕极轴旋转而成的，因此对应于 $0 \leqslant \theta \leqslant \pi$，按照极坐标的弧微分公式，有
\[
\sqrt{r^2 + r'^2} = \sqrt{a^2(1 + \cos\theta)^2 + a^2\sin^2\theta} = a\sqrt{2(1 + \cos\theta)} = 2a\left|\cos\frac{\theta}{2}\right|,
\]
而
\[
y = a(1 + \cos\theta)\sin\theta = 4a\cos^3\frac{\theta}{2}\sin\frac{\theta}{2},
\]
故侧面积
\begin{align*}
S &= 2\pi\int_0^\pi y(\theta)\sqrt{[r(\theta)]^2 + [r'(\theta)]^2}\,\mathrm{d}\theta \\
&= 2\pi\int_0^\pi 4a\cos^3\frac{\theta}{2}\sin\frac{\theta}{2}\cdot 2a\left|\cos\frac{\theta}{2}\right|\mathrm{d}\theta \\
&= 16\pi a^2\int_0^\pi \cos^4\frac{\theta}{2}\sin\frac{\theta}{2}\,\mathrm{d}\theta = \frac{32}{5}\pi a^2.
\end{align*}
\end{proof}


\subsection{平面曲线弧与平面图形的质心}\label{subsec:8-2-7}

假设平面上有 $n$ 个质点：
\[
A_1(x_1, y_1), A_2(x_2, y_2), \cdots, A_n(x_n, y_n),
\]
它们的质量分别为 $m_1, m_2, \cdots, m_n$，那么第 $i$ 个质点对 $x$ 轴、$y$ 轴的静力矩（又叫一阶矩）分别为
\[
m_i y_i, \qquad m_i x_i \quad (i = 1, 2, \cdots, n).
\]
因此质点组对 $x$ 轴、$y$ 轴的静力矩分别为
\[
\sum_{i=1}^n m_i y_i, \quad \sum_{i=1}^n m_i x_i.
\]
设该质点组的质心（即质量中心）在点 $\bar{A}(\bar{x}, \bar{y})$ 处，并且记质点组的质量为 $M = \sum_{i=1}^n m_i$，则由静力矩定律知
\[
M_y = \sum_{i=1}^n m_i x_i = M\bar{x}, \quad M_x = \sum_{i=1}^n m_i y_i = M\bar{y},
\]
从而质点组的质心坐标为
\[
\bar{x} = \frac{M_y}{M} = \frac{\sum_{i=1}^n m_i x_i}{M}, \qquad \bar{y} = \frac{M_x}{M} = \frac{\sum_{i=1}^n m_i y_i}{M}.
\]

现在考虑一条质量均匀分布的平面物质曲线弧 $\wideparen{AB}$，其线密度为常数 $\rho$．我们用微元法来确定 $\wideparen{AB}$ 的质心 $G(\bar{x}, \bar{y})$．

记曲线弧 $\wideparen{AB}$ 的长度为 $l$．取 $A$ 点作为计算弧长的起点，并取弧长 $s$ 为自变量，显然有 $0 \leqslant s \leqslant l$．

分割弧长区间 $[0, l]$，任取一份 $[s, s + \mathrm{d}s]$，我们可以近似把它看作一个质点，其坐标为 $(x, y)$（图~\ref{fig:8-27}），易知这一小段曲线弧的质量为 $\rho\,\mathrm{d}s$，于是其对 $x$ 轴、$y$ 轴的静力矩微元分别为
\[
\mathrm{d}M_x = y\rho\,\mathrm{d}s, \quad \mathrm{d}M_y = x\rho\,\mathrm{d}s.
\]
从 $0$ 到 $l$ 求定积分，得到曲线弧 $\wideparen{AB}$ 对 $x$ 轴和 $y$ 轴的静力矩
\[
M_x = \int_0^l y\rho\,\mathrm{d}s = \rho\int_0^l y\,\mathrm{d}s,
\]
\[
M_y = \int_0^l x\rho\,\mathrm{d}s = \rho\int_0^l x\,\mathrm{d}s.
\]
此外，不难用微元法求得曲线弧 $\wideparen{AB}$ 的质量
\[
M = \int_0^l \rho\,\mathrm{d}s = \rho\int_0^l \mathrm{d}s = \rho l.
\]
于是得到曲线弧 $\wideparen{AB}$ 的质心坐标
\[
\bar{x} = \frac{M_y}{M} = \frac{\rho\int_0^l x\,\mathrm{d}s}{\rho l} = \frac{\int_0^l x\,\mathrm{d}s}{l},
\]
\[
\bar{y} = \frac{M_x}{M} = \frac{\rho\int_0^l y\,\mathrm{d}s}{\rho l} = \frac{\int_0^l y\,\mathrm{d}s}{l}.
\]
上两式的两端同乘 $2\pi l$，得到
\[
2\pi\bar{x}l = 2\pi\int_0^l x\,\mathrm{d}s = F_1,
\]
\[
2\pi\bar{y}l = 2\pi\int_0^l y\,\mathrm{d}s = F_2,
\]
其中 $F_1, F_2$ 分别为曲线弧 $\wideparen{AB}$ 绕 $y$ 轴和 $x$ 轴旋转所得旋转体的侧面积，由此得到

\begin{theorem}[古鲁金第一定理]\label{thm:8-5}
质量均匀分布的平面曲线绕同一平面上某一条不与其相交的轴旋转，由此所产生的旋转体的侧面积 $F$，等于曲线弧的质心绕同一轴旋转所生成的圆周之长乘以该曲线弧的弧长 $l$．
\end{theorem}

这个定理告诉我们：已知弧长 $l$ 与侧面积 $F$，可以求质心 $(\bar{x}, \bar{y})$；已知弧长 $l$ 与质心 $(\bar{x}, \bar{y})$，又可以求侧面积 $F$．

\begin{figure}[htbp]
\centering
\includegraphics[width=0.45\textwidth]{fig8-27.png}
\figcap{图 8-27}\label{fig:8-27}
\end{figure}

\begin{example}
求半径为 $R$ 的半圆周的质心．
\end{example}

\begin{proof}[解]
取坐标系如图~\ref{fig:8-28} 所示．由于半圆周对称于 $y$ 轴，且质量均匀分布，因此质心 $(\bar{x}, \bar{y})$ 必在 $y$ 轴上，即 $\bar{x} = 0$，只需求 $\bar{y}$．显然，半圆周的长度为 $l = \pi R$，半圆周绕 $x$ 轴旋转所产生的旋转体的侧面积为 $F = 4\pi R^2$．由古鲁金第一定理知 $4\pi R^2 = 2\pi\bar{y}\cdot\pi R$，因此
\[
\bar{y} = \frac{2}{\pi}R.
\]
从而得到质心 $\left(0, \frac{2}{\pi}R\right)$．
\end{proof}

\begin{figure}[htbp]
\centering
\begin{minipage}[t]{0.45\textwidth}
\centering
\includegraphics[width=0.85\textwidth]{fig8-28.png}
\figcap{图 8-28}\label{fig:8-28}
\end{minipage}
\hfill
\begin{minipage}[t]{0.45\textwidth}
\centering
\includegraphics[width=0.85\textwidth]{fig8-29.png}
\figcap{图 8-29}\label{fig:8-29}
\end{minipage}
\end{figure}

\begin{example}
求由圆周 $x^2 + (y-a)^2 = R^2$ ($a > R > 0$) 绕 $x$ 轴旋转所成圆环体的侧面积（图~\ref{fig:8-29}）．
\end{example}

\begin{proof}[解]
圆周长 $l = 2\pi R$，圆周的质心坐标为 $\bar{x} = 0, \bar{y} = a$．由古鲁金第一定理得到圆环体的侧面积
\[
F = l\cdot 2\pi\bar{y} = 2\pi R\cdot 2\pi a = 4\pi^2 Ra.
\]
这种解法比直接利用侧面积公式来计算要简单得多．
\end{proof}

类似地我们可求平面图形（图~\ref{fig:8-30}）
\[
y_1(x) \leqslant y \leqslant y_2(x), \quad a \leqslant x \leqslant b
\]
的质心．设图形的密度是均匀的，不妨设 $\rho = 1$．

\begin{figure}[htbp]
\centering
\includegraphics[width=0.45\textwidth]{fig8-30.png}
\figcap{图 8-30}\label{fig:8-30}
\end{figure}

在 $x$ 的附近取面积微元，则它对 $x, y$ 轴的静力矩分别是
\[
\mathrm{d}M_x = \frac{y_2 + y_1}{2}(y_2 - y_1)\,\mathrm{d}x,
\]
\[
\mathrm{d}M_y = x(y_2 - y_1)\,\mathrm{d}x,
\]
因此
\[
M_x = \int_a^b \frac{y_2^2 - y_1^2}{2}\,\mathrm{d}x,
\]
\[
M_y = \int_a^b x(y_2 - y_1)\,\mathrm{d}x.
\]
记平面图形的面积为 $S$，即
\[
S = \int_a^b (y_2 - y_1)\,\mathrm{d}x,
\]
则质心的坐标为
\[
\begin{cases}
\bar{x} = \displaystyle\frac{1}{S}\int_a^b x(y_2 - y_1)\,\mathrm{d}x, \\
\bar{y} = \displaystyle\frac{1}{S}\int_a^b \frac{y_2^2 - y_1^2}{2}\,\mathrm{d}x.
\end{cases}
\]
在后一等式中两边乘以 $2\pi S$，则得
\[
2\pi\bar{y}S = \pi\int_a^b (y_2^2 - y_1^2)\,\mathrm{d}x.
\]
等式右边刚巧等于平面图形绕 $x$ 轴旋转所得旋转体的体积，故有下面的

\begin{theorem}[古鲁金第二定理]\label{thm:8-6}
质量均匀分布的平面图形绕此平面上一条与之不相交的直线旋转，所得旋转体之体积，等于质心旋转时走过的圆周长乘以平面图形的面积．
\end{theorem}

\begin{example}
求半径为 $R$ 的半圆域的质心坐标（参见图~\ref{fig:8-28}）．
\end{example}

\begin{proof}[解]
取圆心为坐标原点，半圆直径为 $x$ 轴，则质心必在 $y$ 轴上，故只需求质心的纵坐标．根据古鲁金第二定理，有
\[
2\pi\bar{y}\cdot\frac{\pi}{2}R^2 = \frac{4}{3}\pi R^3,
\]
于是
\[
\bar{y} = \frac{4}{3\pi}R.
\]
\end{proof}


\subsection{转动惯量}\label{subsec:8-2-8}

从物理学知，质量为 $m$ 的质点绕固定轴旋转时，其转动惯量（即二阶矩）为 $J = mr^2$，其中 $r$ 表示质点到固定轴的距离．

设有 $n$ 个质量分别为 $m_1, m_2, \cdots, m_n$ 的质点，它们到某固定轴的距离分别为 $r_1, r_2, \cdots, r_n$，则此质点组对固定轴的转动惯量为
\[
J = m_1 r_1^2 + m_2 r_2^2 + \cdots + m_n r_n^2 = \sum_{i=1}^n m_i r_i^2.
\]
如果是一个质量连续分布的物体绕固定轴转动，那么，怎样求转动惯量呢？一般说来，需要用到重积分，或曲线积分、曲面积分．但是，当物体的质量均匀分布，且物体的形状具有某种对称性时，有时也可用定积分来解决．

\begin{example}
设有一均匀细杆，长为 $2l$，质量为 $M$，固定轴 $u$ 通过细杆的中心且与细杆垂直（图~\ref{fig:8-31}）．求细杆对轴 $u$ 的转动惯量 $J_u$．
\end{example}

\begin{proof}[解]
取坐标系如图~\ref{fig:8-31} 所示．分割区间 $[-l, l]$，任取一小段 $[x, x + \mathrm{d}x]$．这一小段细杆可近似看成一个质点，它到轴 $u$ 的距离为 $|x|$，质量为 $\frac{M}{2l}\mathrm{d}x$．于是得到转动惯量微元
\[
\mathrm{d}J_u = \left(\frac{M}{2l}\mathrm{d}x\right)x^2 = \frac{M}{2l}x^2\,\mathrm{d}x.
\]
从而转动惯量为
\[
J_u = \int_{-l}^l \frac{M}{2l}x^2\,\mathrm{d}x = \frac{1}{3}Ml^2.
\]
\end{proof}

\begin{figure}[htbp]
\centering
\begin{minipage}[t]{0.45\textwidth}
\centering
\includegraphics[width=0.9\textwidth]{fig8-31.png}
\figcap{图 8-31}\label{fig:8-31}
\end{minipage}
\hfill
\begin{minipage}[t]{0.45\textwidth}
\centering
\includegraphics[width=0.9\textwidth]{fig8-32.png}
\figcap{图 8-32}\label{fig:8-32}
\end{minipage}
\end{figure}

\begin{example}
设有一质量为 $M$、半径为 $R$ 的均匀圆盘，求圆盘对于它的一条直径的转动惯量 $J$．
\end{example}

\begin{proof}[解]
取坐标系如图~\ref{fig:8-32} 所示，$x$ 轴与某条直径重合．分割区间 $[-R, R]$，圆盘相应地被分成若干平行于 $y$ 轴的小窄条．任取一份 $[x, x + \mathrm{d}x]$，相应于它的小窄条可近似看成一个细杆，其长度为 $2y = 2\sqrt{R^2 - x^2}$，质量为
\[
\mathrm{d}m = \text{面密度}\times\text{面积} = \frac{M}{\pi R^2}(2y\,\mathrm{d}x) = \frac{2M}{\pi R^2}\sqrt{R^2 - x^2}\,\mathrm{d}x.
\]
由上例知，此细杆对 $x$ 轴的转动惯量为
\[
\mathrm{d}J = \frac{1}{3}(\mathrm{d}m)y^2 = \frac{2M}{3\pi R^2}(R^2 - x^2)^{3/2}\,\mathrm{d}x.
\]
从而得到所求转动惯量为
\begin{align*}
J &= \int_{-R}^R \frac{2M}{3\pi R^2}(R^2 - x^2)^{3/2}\,\mathrm{d}x \\
&= \frac{4M}{3\pi R^2}\int_0^R (R^2 - x^2)^{3/2}\,\mathrm{d}x \quad (\text{令 } x = R\sin t) \\
&= \frac{4M}{3\pi R^2}\int_0^{\pi/2} R^4\cos^4 t\,\mathrm{d}t = \frac{1}{4}MR^2.
\end{align*}
\end{proof}

\begin{example}
设有一质量为 $M$、半径为 $R$ 的均匀物质圆周，求它对通过其中心且垂直于圆周所在平面的固定轴 $u$ 的转动惯量（图~\ref{fig:8-33}）．
\end{example}

\begin{proof}[解]
分割圆周为若干小弧段，任取一段 $[s, s + \mathrm{d}s]$，它可近似看成一个质点，其上集中了小弧段 $\mathrm{d}s$ 的质量 $\mathrm{d}m$，易知
\[
\mathrm{d}m = \text{线密度}\times\text{长度} = \frac{M}{2\pi R}\mathrm{d}s,
\]
这个质点到轴 $u$ 的距离为 $R$，因此它对轴 $u$ 的转动惯量为
\[
\mathrm{d}J = (\mathrm{d}m)R^2 = \frac{MR}{2\pi}\mathrm{d}s.
\]
于是得到圆周对轴 $u$ 的转动惯量
\[
J = \int_0^{2\pi R} \frac{MR}{2\pi}\mathrm{d}s = \frac{MR}{2\pi}\int_0^{2\pi R}\mathrm{d}s = \frac{MR}{2\pi}\cdot 2\pi R = MR^2.
\]
\end{proof}

\begin{figure}[htbp]
\centering
\begin{minipage}[t]{0.45\textwidth}
\centering
\includegraphics[width=0.75\textwidth]{fig8-33.png}
\figcap{图 8-33}\label{fig:8-33}
\end{minipage}
\hfill
\begin{minipage}[t]{0.45\textwidth}
\centering
\includegraphics[width=0.85\textwidth]{fig8-34.png}
\figcap{图 8-34}\label{fig:8-34}
\end{minipage}
\end{figure}

\begin{example}
设有一质量为 $M$、半径为 $R$ 的均匀圆盘，求它对通过圆心且与盘面垂直的轴 $u$ 的转动惯量．
\end{example}

\begin{proof}[解]
取坐标系如图~\ref{fig:8-34} 所示．分割 $x$ 轴上的区间 $[0, R]$，圆盘相应地被分成若干个窄圆环．任取一个小区间 $[x, x + \mathrm{d}x]$，相应于它的，是一个内半径为 $x$，外半径为 $x + \mathrm{d}x$ 的窄圆环，其质量为
\begin{align*}
\mathrm{d}m &= \text{面密度}\times\text{面积} = \frac{M}{\pi R^2}\left[\pi(x + \mathrm{d}x)^2 - \pi x^2\right] \\
&= \frac{M}{\pi R^2}\left[2\pi x\,\mathrm{d}x + \pi(\mathrm{d}x)^2\right].
\end{align*}
由于 $\mathrm{d}x$ 很小，因此可忽略 $(\mathrm{d}x)^2$，于是有质量微元
\[
\mathrm{d}m = \frac{2M}{R^2}x\,\mathrm{d}x.
\]
这个窄圆环可近似看成半径为 $x$ 的圆周，于是由上例知，它对 $u$ 轴的转动惯量为
\[
\mathrm{d}J = (\mathrm{d}m)x^2 = \frac{2M}{R^2}x^3\,\mathrm{d}x,
\]
从而得到圆盘对 $u$ 轴的转动惯量
\[
J = \int_0^R \frac{2M}{R^2}x^3\,\mathrm{d}x = \frac{1}{2}MR^2.
\]
\end{proof}

\begin{figure}[htbp]
\centering
\includegraphics[width=0.45\textwidth]{fig8-35.png}
\figcap{图 8-35}\label{fig:8-35}
\end{figure}

\begin{example}
设有一质量为 $M$、半径为 $R$ 的均匀球体，求它对其直径的转动惯量．
\end{example}

\begin{proof}[解]
取坐标系如图~\ref{fig:8-35} 所示，$x$ 轴通过球的直径．分割 $x$ 轴上的区间 $[-R, R]$，球体相应地被分成若干个形状为球台的薄片．任取一个小区间 $[x, x + \mathrm{d}x]$，相应于它的薄球台可近似看成一个半径为 $y = \sqrt{R^2 - x^2}$ 的薄圆盘，其厚度为 $\mathrm{d}x$．易知此圆盘的体积为 $\pi y^2\,\mathrm{d}x = \pi(R^2 - x^2)\,\mathrm{d}x$，从而其质量为
\[
\mathrm{d}m = \text{体密度}\times\text{体积} = \frac{M}{\frac{4}{3}\pi R^3}\cdot\pi(R^2 - x^2)\,\mathrm{d}x = \frac{3M}{4R^3}(R^2 - x^2)\,\mathrm{d}x.
\]
因为 $\mathrm{d}x$ 很小，所以该薄圆盘可按上例计算转动惯量，得到
\[
\mathrm{d}J = \frac{1}{2}(\mathrm{d}m)y^2 = \frac{3M}{8R^3}(R^2 - x^2)^2\,\mathrm{d}x,
\]
从而球体对直径的转动惯量为
\[
J = \int_{-R}^R \frac{3M}{8R^3}(R^2 - x^2)^2\,\mathrm{d}x = \frac{2}{5}MR^2.
\]
\end{proof}


\begin{xiti}
\item 求下列各曲线所围成的图形面积：
\begin{ziti}
\item $y^2 = 4(x+1)$，$y^2 = 4(1-x)$；
\item $y = |\ln x|$，$y = 0$ ($0.1 \leqslant x \leqslant 10$)；
\item $y = x$，$y = x + \sin^2 x$ ($0 \leqslant x \leqslant \pi$)；
\item $y^2 = 2x$，$x = 5$；
\item $y = x^2$，$y = x + 5$；
\item $x^{2/3} + y^{2/3} = a^{2/3}$．
\end{ziti}

\item 求下列用极坐标表示的曲线所围图形的面积：
\begin{ziti}
\item 双扭线 $r^2 = a^2\cos 2\varphi$；
\item 三叶玫瑰线 $r = a\sin 3\varphi$；
\item 蚌线 $r = a\cos\theta + b$ ($b \geqslant a$).
\end{ziti}

\item 求下列用参数方程表示的曲线所围图形的面积：
\begin{ziti}
\item $x = 2t - t^2$，$y = 2t^2 - t^3$；
\item 摆线 $x = a(t - \sin t)$，$y = a(1 - \cos t)$ ($0 \leqslant t \leqslant 2\pi$) 及 $x$ 轴；
\item 圆的渐开线 $x = a(\cos t + t\sin t)$，$y = a(\sin t - t\cos t)$ ($0 \leqslant t \leqslant 2\pi$) 及半直线 $x = a$ ($y \leqslant 0$)，其中 $a > 0$．
\end{ziti}

\item 直线 $y = x$ 把椭圆 $x^2 + 3y^2 = 6y$ 的面积分成两部分 $A$（小的一块）和 $B$（大的一块），求 $\frac{A}{B}$ 之值．

\item 求 $r = 3\cos\theta$ 和 $r = 1 + \cos\theta$ 所围的公共部分的面积．

\item 求下列旋转体的体积：
\begin{ziti}
\item 椭圆 $\frac{x^2}{a^2} + \frac{y^2}{b^2} = 1$ 绕 $x$ 轴；
\item $y = \sin x$，$y = 0$ ($0 \leqslant x \leqslant \pi$)：(i) 绕 $x$ 轴，(ii) 绕 $y$ 轴；
\item 旋轮线 $x = a(t - \sin t)$，$y = a(1 - \cos t)$ ($0 \leqslant t \leqslant 2\pi$), $y = 0$：(i) 绕 $x$ 轴，(ii) 绕 $y$ 轴，(iii) 绕直线 $y = 2a$；
\item 双曲线 $\frac{y^2}{b^2} - \frac{x^2}{a^2} = 1$ 与直线 $x = \pm h$ 所围的图形绕 $x$ 轴旋转．
\end{ziti}

\item 求由下列各曲面所围成的几何体的体积：
\begin{ziti}
\item 求截锥体的体积，其上、下底皆为椭圆，椭圆的轴长分别等于 $A, B$ 和 $a, b$，而高为 $h$；
\item 正圆台：其上、下底分别是半径为 $a, b$ 的圆，而其间的距离为 $h$．
\end{ziti}

\item 已知球半径为 $R$，试求高为 $h$ 的球冠体积 ($h \leqslant R$)．

\item 求下列曲线的弧长：
\begin{ziti}
\item $y = x^2$，$0 \leqslant x \leqslant 1$；
\item $y = \mathrm{e}^x$，$1 \leqslant x \leqslant 2$；
\item $\sqrt{x} + \sqrt{y} = 1$；
\item 星形线 $x = a\cos^3 t$，$y = a\sin^3 t$ ($0 \leqslant t \leqslant 2\pi$)；
\item 圆的渐开线 $x = a(\cos t + t\sin t)$，$y = a(\sin t - t\cos t)$，$a > 0, 0 \leqslant t \leqslant 2\pi$；
\item $r = a\sin^3\frac{\theta}{3} \quad (a > 0)$；
\item 心脏线 $r = a(1 + \cos\theta), 0 \leqslant \theta \leqslant 2\pi, a > 0$．
\end{ziti}

\item 求下列各曲线在指定点的曲率和曲率半径：
\begin{ziti}
\item $xy = 4$ 在点 $(2, 2)$；
\item $y = \ln x$ 在点 $(1, 0)$．
\end{ziti}

\item 求下列曲线的曲率与曲率半径：
\begin{ziti}
\item 抛物线 $y^2 = 2px$ ($p > 0$)；
\item 双曲线 $\frac{x^2}{a^2} - \frac{y^2}{b^2} = 1$；
\item 星形线 $x^{2/3} + y^{2/3} = a^{2/3}$．
\end{ziti}

\item 求下列参数方程给出的曲线的曲率和曲率半径：
\begin{ziti}
\item 旋轮线 $x = a(t - \sin t)$，$y = a(1 - \cos t)$ ($a > 0$)；
\item 椭圆 $x = a\cos t$，$y = b\sin t$ ($a, b > 0$)；
\item 圆的渐开线 $x = a(\cos t + t\sin t)$，$y = a(\sin t - t\cos t)$．
\end{ziti}

\item 求下列以极坐标表示的曲线的曲率半径：
\begin{ziti}
\item 心脏线 $r = a(1 + \cos\theta)$ ($a > 0$)；
\item 双扭线 $r^2 = 2a^2\cos 2\theta$ ($a > 0$)；
\item 对数螺线 $r = a\mathrm{e}^{\lambda\theta}$ ($\lambda > 0$).
\end{ziti}

\item 设曲线是用极坐标方程 $r = r(\theta)$ 给出，且二阶可导，证明它在点 $\theta$ 处曲率为
\[
K = \frac{|r^2 + 2r'^2 - r r''|}{(r^2 + r'^2)^{3/2}}.
\]

\item 证明抛物线 $y = ax^2 + bx + c$ 在顶点处的曲率半径为最小．

\item 求曲线 $y = 2(x - 1)^2$ 的最小曲率半径．

\item 求曲线 $y = \mathrm{e}^x$ 上曲率最大的点．

\item 求下列平面曲线绕轴旋转所得旋转曲面的面积：
\begin{ziti}
\item $y = \sin x$，$0 \leqslant x \leqslant \pi$ 绕 $x$ 轴；
\item $x = a(t - \sin t)$，$y = a(1 - \cos t)$，$a > 0$，$0 \leqslant t \leqslant 2\pi$ 绕直线 $y = 2a$；
\item $\frac{x^2}{a^2} + \frac{y^2}{b^2} = 1$ ($a > b$) 绕 $x$ 轴；
\item $x = a\cos^3 t$，$y = a\sin^3 t$ 绕 $x$ 轴；
\item $r^2 = 2a^2\cos 2\theta$ 绕极轴．
\end{ziti}

\item 求下列曲线段的质心：
\begin{ziti}
\item 半径为 $r$，弧长为 $\frac{1}{2}\pi\alpha$ ($\alpha \leqslant \pi$) 的均匀圆弧；
\item 对数螺线 $r = a\mathrm{e}^{k\theta}$ ($a > 0, k > 0$) 上由点 $(0, a)$ 到点 $(\theta, r)$ 的均匀弧段；
\item 以 $A(0,0), B(0,1), C(2,1), D(2,0)$ 为顶点的矩形周界，曲线上任一点的密度等于该点到原点距离的 $2$ 倍；
\item $x = a(t - \sin t)$，$y = a(1 - \cos t)$，$0 \leqslant t \leqslant 2\pi$，$a > 0$，密度为常数．
\end{ziti}

\item 已知一抛物线段 $y = x^2$ ($-1 \leqslant x \leqslant 1$)，曲线段上任一点处的密度与该点到 $y$ 轴的距离成正比，$x = 1$ 处密度为 $5$，求此曲线段的质量．

\item 轴长 $10\text{ m}$，密度分布为 $\rho(x) = (6 + 0.3x)\text{ kg/m}$，其中 $x$ 为距轴的一个端点的距离，求轴的质量．

\item 求半球 $0 \leqslant z \leqslant \sqrt{R^2 - x^2 - y^2}$ 的质心．

\item 求锥体 $\sqrt{x^2 + y^2} \leqslant z \leqslant h$ 的质心和绕 $z$ 轴的转动惯量．

\item 求抛物体 $x^2 + y^2 \leqslant z \leqslant h$ 的质心和绕 $z$ 轴的转动惯量．
\end{xiti}
